% Source-aligned English for printed p.134.
\SourcePage{134}

\LayerTarget{tohoku:prop:1.8:proof:p2}
\noindent $=w_{ij}(A_i)$, let $N$ be the subobject of $S$ that is the supremum
of the $N_{ij}$ (which exists by axiom AB 3)), put $\overline A=S/N$, and let
$u_i:A_i\longrightarrow\overline A$ be the morphism induced by the canonical
morphism $S\longrightarrow S/N$. One sees immediately that
$(\overline A,u_i)$ is an inductive limit of $\mathbf A$. We leave to the
reader the proof of the other assertions of Proposition 1.8, a proof that is
of course well known.

\LayerTarget{tohoku:section:1.9}
\noindent\textbf{1.9. Generators and cogenerators.} Let $C$ be a category and
let $(U_i)_{i\in I}$ be a family of objects of $C$. It is called a
\textit{family of generators of $C$} if, for every object $A\in C$ and every
proper subobject $B\neq A$, there are an $i\in I$ and a morphism
$u:U_i\longrightarrow A$ that does not come from a morphism from $U_i$ to
$B$. Then, for every $A\in C$, the subobjects of $A$ form a \textit{set}:
indeed, a subobject $B$ of $A$ is completely determined by the set of
morphisms from objects $U_i$ to $A$ that come from a morphism from $U_i$ to
$B$. An object $U\in C$ is called a \textit{generator} of $C$ if the singleton
family $\{U\}$ is a family of generators.

\LayerTarget{tohoku:prop:1.9.1}
\noindent\textsc{Proposition 1.9.1.} --- \textit{Suppose that $C$ is an
abelian category satisfying axiom AB 3), the existence of infinite direct
sums. Let $(U_i)_{i\in I}$ be a family of objects of $C$, and let
$U=\bigoplus U_i$ be its direct sum. The following conditions are equivalent:}

\begin{enumerate}
\item[a)] \textit{$(U_i)$ is a family of generators of $C$.}
\item[b)] \textit{$U$ is a generator of $C$.}
\item[c)] \textit{Every $A\in C$ is isomorphic to a quotient of a direct sum
$U^{(I)}$ of objects all identical to $U$.}
\end{enumerate}

\LayerTarget{tohoku:prop:1.9.1:proof}
The equivalence of a) and b) follows almost immediately from the definition.
Condition b) implies c): take $I=\mathrm{Hom}(U,A)$ and consider the morphism
from $U^{(I)}$ to $A$ whose component indexed by each $u\in I$ is $u$ itself.
The image $B$ of this morphism is $A$, since otherwise there would be a
$u\in\mathrm{Hom}(U,A)=I$ such that $u(U)\not\subset B$, which is absurd.
Thus $A$ is isomorphic to a quotient of $U^{(I)}$. Condition c) implies b),
because if $A$ is a quotient of $U^{(I)}$, then for every proper subobject
$B$ of $A$ there is an $i$ such that the canonical image in $A$ of the
$i$-th factor of $U^{(I)}$ is not contained in $B$; this gives a morphism from
$U$ to $A$ that does not come from a morphism from $U$ to $B$. (Notice that
the additive structure of $C$ has not really been used here.)

\LayerTarget{tohoku:section:1.9:examples:p1}
\noindent\textsc{Examples.} If $C$ is the abelian category of unital left
modules over a ring $U$ with identity, then $U$, regarded as a left module over
itself, is a generator. If $C$ is the category of sheaves of abelian groups on
a fixed topological space $X$, and for every open $U\subset X$ we denote by
$Z_U$ the sheaf on $X$ that is zero over $X\setminus U$ and equal to the
constant sheaf of integers over $U$, then the family of the $Z_U$ forms a
system of generators of $C$. This example immediately generalizes to the case
of a sheaf $O$ of rings on $X$ and the category of sheaves of $O$-modules on
$X$. Further examples are contained in the following proposition:

\TohokuContinuation{135}
